\section{Related Work}
\teraheapName{} combines techniques from several areas, including memory
management and storage. Thus, we group the related work in the
following categories:
%
(1) scaling managed heaps beyond local DRAM,
%
(2) region-based memory management, and
%
(3) mitigating \teraheapSD overhead.

\subsection{Taxonomy of Scaling Managed Heaps Beyond DRAM}
Recent efforts target local block-addressable storage devices, such as
hard disk drives (HDD) and NVMe SSDs, and local
byte-addressable NVM, or remote memory (DRAM) for storing managed
heaps beyond DRAM.

\paragraph{\textbf{Scaling managed heaps with local block-addressable
storage devices}}
One line of recent work focuses on using block-addressable storage
devices. Melt~\cite{bond:melt} and Leaksurvivor~\cite{leaksurvivor}
use a primary heap over DRAM and a second heap over an HDD for moving
cold objects that programs do not use. Unlike \teraheapName{}, these works do
not use a memory-mapped I/O but maintain a specialized lookup
mechanism in existing JVM code to track the offset of each object in
the device. Although this technique targets the high latency of the
hard disk drives, it requires multiple systems calls and incurs the
high hit overhead of user space lookups in every load/store operation.
Also, LeakSurvivor~\cite{leaksurvivor} leverages the capacity of the
HDD to perform lazy object reclamation. However, it reclaims objects
in the HDD only when the device becomes full, paying high \teraheapGC cost for
long-running applications. Performing \teraheapGC scans over a storage device
introduces high I/O traffic, increasing \teraheapGC pause time.
TMO~\cite{weiner:tmo} monitors application DRAM usage and
transparently offloads cold data to NVMe SSD. Unlike these works,
\teraheapName{} controls which objects to move to the second heap and
eliminates slow \teraheapGC traversals over objects on NVMe SSD. Finally,
prior effort~\cite{kolokasis:teracache} discusses a preliminary
prototype of \teraheapName{}~\cite{teraheap}, while this article presents the
design and implementation of \teraheapName{} in OpenJDK17.

\paragraph{\textbf{Scaling managed heaps with byte-addressable devices}}
Recent research efforts focus on extending the managed heaps with
local byte-addressable NVM or remote DRAM. Akram et
al.~\cite{akram:ration, akram:crystal}, GCMove~\cite{gcmove}, and
ThinGC~\cite{thingc} focus on improving NVM write endurance. They
select which objects to move from DRAM to NVM based on their hotness.
Existing work uses two techniques to track object hotness: (1) through
load reference barriers used by concurrent copying collectors to
identify the most accessed objects, and (2) through post-write
reference barriers to identify the most written/updated objects. Both
techniques increase the cost of the barriers and can work with
specific garbage collectors. NVM-friendly-GC~\cite{Xu:gc_evaluation}
report high \teraheapGC overhead with NVM-backed volatile heaps and optimize
the G1 \teraheapGC for Intel Optane persistent memory. It allocates its heap
entirely on NVM and uses DRAM as a user-space cache for \teraheapGC
operations. However, all the accesses to objects by application
(mutator) threads are served by NVM. Panthera~\cite{wang:Panthera}
extends the managed heap over hybrid DRAM and non-volatile memory
(NVM) to scale on-heap caching in Spark. Panthera increases \teraheapGC
overhead as scanning and compacting objects on the managed NVM heap
costs more than collecting the DRAM heap. Also,
JDPHeap~\cite{jpd_heap} uses a managed heap in DRAM and a second
managed heap over NVM. It introduces an extensive programming model
that enables users to select which objects to move from DRAM to NVM
and vice versa. However, unlike \teraheapName{}, JPDHeap cannot avoid
expensive \teraheapGC scans over objects in NVM.

Also, there is an amount of effort~\cite{wang:memliner, wang:semeru,
ma:mako} for extending the managed heaps using remote memory. On the
one hand, Memliner~\cite{wang:memliner} is designed for
software-disaggregated datacenters. MemLiner reorganizes the access
order of the objects by the \teraheapGC threads to follow a similar (not
identical) memory-access path with the mutator threads to enhance the
speed of far-memory accesses during \teraheapGC. On the other hand,
Semeru~\cite{wang:semeru} and Mako~\cite{ma:mako} target
hardware-designed datacenters. Semeru supports running managed big
data frameworks on disaggregated hardware by dividing the conventional
JVM into two instances:
%
(1) the CPU-JVM instance that executes the program on the CPU server
and
%
(2) the memory-JVM instance that performs \teraheapGC on the memory
server.
%
To maximize \teraheapGC performance, Semeru proposes a distributed garbage
collector that offloads the liveness analysis for Java objects to
memory servers concurrently with mutator threads. Unlike Semeru, Mako
offloads to memory servers both object scanning for liveness analysis
and evacuation, reducing \teraheapGC pause time. Although DRAM is
byte-addressable, these approaches rely on the OS-swap system to move
objects between local and remote memory at page granularity. \teraheapName{}
is orthogonal to these approaches because it exploits fast local storage devices for scaling the
capacity of the managed heap.

\setlength{\unitlength}{1em}

\begin{table}[t]
\centering
\caption{Comparison of selected systems that extend managed heaps beyond DRAM.}
\resizebox{0.8\columnwidth}{!}{%
\begin{tabular}{lccccc}
\hline
\multicolumn{1}{c}{\multirow{2}{*}{\textbf{System}}} &
  \multirow{2}{*}{\textbf{\begin{tabular}[c]{@{}c@{}}Developer \\ Effort\end{tabular}}} &
  \multicolumn{2}{c}{\textbf{\begin{tabular}[c]{@{}c@{}}GC Scans\\ on the Device\end{tabular}}} &
  \multicolumn{2}{c}{\textbf{Device Type}} \\ \cline{3-6}
\multicolumn{1}{c}{}         &      & \textbf{Lazy} & \textbf{Eager}       & \textbf{Block-addressable} & \textbf{Byte-addressable} \\ \hline
  LeakSurvivor~\cite{leaksurvivor}                  & High & (\checkmark)  & -                    & \checkmark                 & -                         \\
  Melt~\cite{bond:melt}                             & None & (\checkmark)  & -                    & \checkmark                 & -                         \\
  Write Rationing~\cite{akram:ration}               & None & -             & \checkmark           & -                          & \checkmark                \\
  Panthera~\cite{wang:Panthera}                     & Low  & -             & \checkmark           & -                          & \checkmark                \\
  NVM-friendly-GC~\cite{nvm-gc-eurosys}             & None & -             & \checkmark           & -                          & \checkmark                \\
  JPD Heap~\cite{jpd_heap}                     & High & -             & \checkmark           & -                          & \checkmark                \\
  GCMove~\cite{gcmove}                      & None & -             & \checkmark           & -                          & \checkmark                \\
  ThinGC~\cite{thingc}                      & None & -             & \checkmark           & -                          & \checkmark                \\
  TMO~\cite{weiner:tmo}                          & None & -             & \checkmark           & \checkmark                 & -                         \\
  Semeru~\cite{wang:semeru}                       & None & -             & \checkmark           & -                          & \checkmark                \\
  MemLiner~\cite{wang:memliner}                     & None & -             & \checkmark           & -                          & \checkmark                \\
  Mako~\cite{ma:mako}                         & None & -             & \checkmark           & -                          & \checkmark                \\ \hline
\textbf{TeraHeap (this thesis)} & Low  & \checkmark    & -                    & \checkmark                 & \checkmark                \\ \hline
\end{tabular}%
}
\label{tab:related}
\end{table}


Table~\ref{tab:related} illustrates the taxonomy of existing
approaches that scale the managed heaps beyond local DRAM capacity
based on the following characteristics:
%
(1) development effort,
%
(2) object reclamation on the device or remote memory and
%
(3) device type.
%
\teraheapName{} with low development effort uses the framework's knowledge
regarding which objects are suitable for moving to the storage device.
Also, \teraheapName{} avoids \teraheapGC scans over the storage device because it leverages the
high capacity of the devices to perform lazy object reclamation.
Finally, it can still improve application performance even if H2 is in
a local block-addressable storage device or a byte-addressable device.

\subsection{Region-based Memory Management}
Managed big data frameworks have started to use region-based memory
management for large heaps. Systems, such as
Gerenuk~\cite{navasca:gerenuk} and Facade~\cite{nguyen:facade},
provide a compilation framework that transforms programmer-specified
classes for off-heap allocation at compile time, thereby reducing the
serialization overhead. However, these systems increase programmer
effort by requiring the developer to specify when to free objects from
native memory and modify legacy applications. Broom~\cite{gog:broom}
uses region annotations but requires refactoring of applications'
source code. Yak~\cite{khanh:yak} requires programmers to annotate
epochs in applications. Yak allocates all objects in an epoch on a
second region-based heap to reduce \teraheapGC time. The epoch abstraction is
appropriate for the map-reduce programming pattern. However, it cannot
handle objects computed lazily or accessed from arbitrary program
locations. Deca~\cite{lu:deca} proposes lifetime-based memory
management for Spark. However, their work only applies to Spark and
cannot be used for other frameworks. Unlike prior work, \teraheapName{}
requires adding hints only in the framework layer. Then, \teraheapName{}
dynamically selects all appropriate objects in the transitive closure
of root objects. NG2C~\cite{bruno:n2gc} uses runtime profiling to
identify long-lived objects. They incur online profiling overhead.
Other work uses offline allocation site profiling to manage
objects~\cite{pretenuring,Blackburn:2001:PJ}. Lifetime profiling can
complement \teraheapName{} and further improve efficiency.

\subsection{Mitigating \teraheapSD Overhead}
Several libraries~\cite{thrift,kryo,protocol-buffers} improve the
efficiency of \teraheapSD, but they cannot reduce high \teraheapGC cost in big data
frameworks. Skyway~\cite{nguyen:skyway} reduces the \teraheapSD cost by
directly transferring objects through the network in distributed
managed heaps, but it does not cope with DRAM limitations and \teraheapGC
overheads. SSDStreamer~\cite{ssdstreamer} is a userspace SSD-based
caching system that uses DRAM as a stream buffer for SSD devices.
Although SSDStreamer reduces \teraheapSD cost by providing a lightweight
serializer, it cannot reduce \teraheapGC cost and the memory pressure in the
managed heap. Recent work~\cite{jang:arch, pourhabibi:optimus}
proposes optimizations for \teraheapSD with custom hardware extensions and
modifications to the programming model. Other
work~\cite{matei:champions, taranov:naos, ibanez:zerializer} focuses
on reducing \teraheapSD cost by reducing the number of object copies across
buffers. This body of work does not mitigate directly \teraheapGC overhead.
\teraheapName{} is the first work that eliminates \teraheapSD cost without increasing
\teraheapGC overhead for a large portion of objects in big data analytics
frameworks.

